% Generated draft: requires mathematical and source review.
\documentclass[11pt]{article}
\usepackage[margin=25mm]{geometry}
\usepackage{fontspec}
\setmainfont{Latin Modern Roman}
\usepackage{amsmath,amssymb,mathtools,unicode-math}
\setmathfont{Latin Modern Math}
\usepackage{enumitem,needspace,xurl,hyperref}
\setlength{\parindent}{0pt}
\setlength{\parskip}{5pt}
\setlength{\emergencystretch}{3em}
\newcommand{\N}{\mathbb{N}}
\newcommand{\Z}{\mathbb{Z}}
\newcommand{\Q}{\mathbb{Q}}
\newcommand{\R}{\mathbb{R}}
\newcommand{\C}{\mathbb{C}}
\newcommand{\F}{\mathbb{F}}
\newcommand{\A}{\mathbb{A}}
\newcommand{\PP}{\mathbb{P}}
\newcommand{\E}{\mathbb{E}}
\newcommand{\eps}{\varepsilon}
\newcommand{\dd}{\,\mathrm{d}}
\begin{document}
\section*{UnsolvedMath ID: 5300041}
\textbf{Source problem code:} catalogue entry
\subsection{Definitions and mathematical statement}
Let $R:\widehat\C\to\widehat\C$ be a rational map of degree at least two. Its Julia set $J(R)$ is the complement of the open set on which the iterates $(R^n)_{n\geq0}$ form a normal family. Thus $J(R)$ is closed. It is nowhere dense exactly when its interior is empty. Let $m_{\mathrm{sph}}$ be spherical area measure; positivity agrees with positivity of planar Lebesgue area in a finite coordinate chart.

The question asks whether
\[
\exists R\text{ rational},\quad \deg R\geq2,\qquad
\operatorname{int}J(R)=\varnothing,\qquad m_{\mathrm{sph}}(J(R))>0.
\]
The supplied solution report says this existence holds even for quadratic polynomials $R(z)=z^2+c$.

Positive area must not be replaced by containing an open region or having full area in an open region. A closed nowhere-dense set cannot have full measure in a nonempty open region, since its complement contains a nonempty open subset there. The supplied summary's parenthetical ``full on a region's area'' and its description via a holomorphic motion of a positive-measure Cantor set are not part of the existence statement and are not justified by the supplied citation description.
\subsection{Short English statement}
Can a Julia set have positive area while containing no open subset? The supplied literature report answers yes, even for quadratic polynomials.
\subsection{Sources}
\begin{itemize}
\item Local UnsolvedMath registry, permanent ID 5300041, exact source code catalogue entry, ``Positive-area nowhere-dense Julia sets''; proposed by Mikhail Lyubich, 1992. Registry status solved, classification SOLVED-IN-LITERATURE.
\item \emph{Stony Brook---Open Problems in Dynamical Systems}, 1992, Lyubich P1, \url{https://www.math.stonybrook.edu/open-problems-dynamical-systems}. Supplied source-TeX and pdf-text extraction accessed July 29, 2026; historical-source provenance presents the item as open and records no comprehensive modern resolution check. No PDF was accessed here.
\item Buff--Ch\'eritat, ``Quadratic Julia sets with positive area'', \emph{Annals of Mathematics} (2012), supplied as the affirmative solution reference. The summary dates the construction programme 2002--2012 and mentions unnamed higher-degree generalizations. Its additional full-area/Cantor-motion descriptions require checking and are not adopted.
\item Provenance flags status and rights NEEDS\_REVIEW: public preprint access and source provenance were verified, but redistribution permission was not established. The later solution is retained as a supplied literature assessment, not a newly performed check.
\end{itemize}
\end{document}
